\documentclass[a4paper,UTF8]{article}
\usepackage{ctex}
\usepackage[margin=1.25in]{geometry}
\usepackage{ifpdf}
\usepackage{color}
\usepackage{graphicx}
\usepackage{amssymb}
\usepackage{amsmath}
\usepackage{amsthm}
\usepackage{enumerate}
\usepackage{bm}
\usepackage{hyperref}
\usepackage{pgfplots}
\usepackage{epsfig}
\usepackage{color}
\usepackage{tcolorbox}
\usepackage{mdframed}
\usepackage{lipsum}
\usepackage{framed}
\usepackage{setspace}
\usepackage{caption}
\usepackage{amsmath}
\usepackage{cases}

\newmdtheoremenv{thm-box}{myThm}
\newmdtheoremenv{prop-box}{Proposition}
\newmdtheoremenv{def-box}{定义}

\setlength{\evensidemargin}{.25in}
\setlength{\textwidth}{6in}
\setlength{\topmargin}{-0.5in}
\setlength{\topmargin}{-0.5in}
% \setlength{\textheight}{9.5in}
%%%%%%%%%%%%%%%%%%此处用于设置页眉页脚%%%%%%%%%%%%%%%%%%
\pgfplotsset{compat=1.18}
\usepackage{fancyhdr}                                
\usepackage{lastpage}                                           
\usepackage{layout}                                             
\footskip = 12pt 
\pagestyle{fancy}                    % 设置页眉                 
\lhead{2024年秋季}                    
\chead{数字信号处理}                                                
% \rhead{第\thepage/\pageref{LastPage}页} 
\rhead{作业三}                                                                                               
\cfoot{\thepage}                                                
\renewcommand{\headrulewidth}{1pt}  			%页眉线宽，设为0可以去页眉线
\setlength{\skip\footins}{0.5cm}    			%脚注与正文的距离           
\renewcommand{\footrulewidth}{0pt}  			%页脚线宽，设为0可以去页脚线

\makeatletter 									%设置双线页眉                                        
\def\headrule{{\if@fancyplain\let\headrulewidth\plainheadrulewidth\fi%
\hrule\@height 1.0pt \@width\headwidth\vskip1pt	%上面线为1pt粗  
\hrule\@height 0.5pt\@width\headwidth  			%下面0.5pt粗            
\vskip-2\headrulewidth\vskip-1pt}      			%两条线的距离1pt        
 \vspace{6mm}}     								%双线与下面正文之间的垂直间距              
\makeatother  

%%%%%%%%%%%%%%%%%%%%%%%%%%%%%%%%%%%%%%%%%%%%%%
\numberwithin{equation}{section}
%\usepackage[thmmarks, amsmath, thref]{ntheorem}
\newtheorem{myThm}{myThm}
\newtheorem*{myDef}{Definition}
\newtheorem*{mySol}{Solution}
\newtheorem*{myProof}{Proof}
\newtheorem*{myRemark}{备注}
\renewcommand{\tilde}{\widetilde}
\renewcommand{\hat}{\widehat}
\newcommand{\indep}{\rotatebox[origin=c]{90}{$\models$}}
\newcommand*\diff{\mathop{}\!\mathrm{d}}

\usepackage{multirow}

%--

%--
\begin{document}

\title{数字信号处理\\
    作业三}
\author{田永铭\, 221900180}
\maketitle

\section*{作业提交注意事项}
\begin{tcolorbox}
    \begin{enumerate}
       \item[(1)] 本次作业提交截止时间为~\textcolor{red}{\textbf{2024/12/31  23:59:59}}，截止时间后不再接收作业；
        \item[(2)] 作业提交方式：使用此~LaTex~模板书写解答，不允许使用手写图片替代~LaTex~格式解题过程，只需提交编译生成的~pdf~文件，将~pdf~文件发送至邮箱：779652674@qq.com；
        \item[(3)] pdf 文件命名方式：学号-姓名-作业号-v版本号, 例~ 1900000-张三-1-v1；如果需要更改已提交的解答，请在截止时间之前提交新版本的解答，并将版本号加一；
        \item[(4)] 作业提交规范占5分，未按照要求提交作业，或~pdf~命名方式不正确，将会被扣除此部分分数。
    \end{enumerate}
\end{tcolorbox}

\newpage
\section{[40pts]}
确定下列离散时间周期信号的傅里叶级数，并写出每一组系数$a_{k}$的模和相位
\begin{enumerate}[(1)]
    \item $x[n]=\sum\limits_{m=-\infty}^{\infty}\{2\delta{[n-4m]}+4\delta{[n-1-4m]}\}$
    \item $x[n]=1-\sin{(\pi n/4)},0\le n\le 11$,且$x[n]$以12为周期
\end{enumerate}

\begin{framed}
    解：{\color{red} 注意，这里的公式跟书本教材保持一致，周期N是除在正变换中的；PPT上是除在逆变换中；两者没有本质区别，请注意批改，如果要除以在逆变换中，$a_n$答案乘一个周期就行,对于傅里叶级数最终表达式是一样的。}
    \begin{enumerate}[(1)]
    	\item 该信号周期为4，系数为：
    	\begin{align*}
    		a_k &= \frac{1}{4}\sum\limits_{n = 0}^{3}\sum\limits_{m=-\infty}^{\infty}\{2\delta{[n-4m]}+4\delta{[n-1-4m]}\}e^{-jnwk}\\
    		&= \frac{1}{4}[2+4e^{-jwk}]\\
    		&= \frac{1}{2} + e^{-j\frac{\pi k}{2}}\quad(k = 0,1,2,3)
    	\end{align*}
    	
    	亦可写为$\frac{1}{2} + cos(\frac{\pi k}{2})-jsin(\frac{\pi k}{2})$,因此有：
    	
    	傅里叶级数为：
    	$x[n] = \sum\limits_{k=0}^{3}a_ke^{j\frac{n\pi k}{2}}$, 其中$a_k = \frac{1}{2} + e^{-j\frac{\pi k}{2}}(k = 0,1,2,3)$.
    	
    	系数的模和相位如下：
    	\begin{itemize}
    		\item \( |a_0| = 1.5, \quad \arg(a_0) = 0 \)
    		\item \( |a_1| = \sqrt{0.5^2 + (-1)^2} = \sqrt{1.25} , \quad \arg(a_1) = -\arctan2 \)
    		\item \( |a_2| = 0.5, \quad \arg(a_2) = \pi \)(负实数相位)
    		\item \( |a_3| = \sqrt{0.5^2 + 1^2} = \sqrt{1.25}, \quad \arg(a_3) = \arctan2 \)
    	\end{itemize}
    	
    	\item  该信号周期为12，系数为：
    	\begin{align*}
    		a_k &= \frac{1}{12}\sum\limits_{n = 0}^{11}(1-sin(\pi n/4))e^{\frac{-jn\pi k}{6}}\\
    		& = \frac{1}{12}\sum\limits_{n = 0}^{11}e^{\frac{-jn\pi k}{6}} - \frac{1}{12}\sum\limits_{n = 0}^{11}sin(\pi n/4)e^{\frac{-jn\pi k}{6}}\quad(k = 0,1,2,\cdots,11)
    	\end{align*}
    
    	\[ x[n] = \sum_{k=<N>} a_k e^{j \frac{2\pi}{N} kn} = \sum_{k=<12>} a_k e^{j \frac{\pi}{6} kn} \]
    	
    	这题将三角化为指数反而更难算，下面一行是通过暴力化为三角函数求解再合并得到的结果，对于每一个$k$ 都对$n$从$0$到$11$暴力求解，得：
    	
    	\[ a_k = \frac{1}{12} \left[ 1 + 2\left(1 - \frac{\sqrt{2}}{2}\right) \cos\left(\frac{k\pi}{6}\right) + 2\left(1 - \frac{\sqrt{2}}{2}\right) \cos\left(\frac{k\pi}{2}\right) \right. \]
    	
    	\[ \left. + 2 \cos\left(\frac{2k\pi}{3}\right) + 2\left(1 + \frac{\sqrt{2}}{2}\right) \cos\left(\frac{5k\pi}{6}\right) + 2(-1)^k \right] \]
    	
    	因此离散傅里叶级数为：
    	$x[n] = \sum_{k=0}^{11}a_ke^{j\frac{n\pi}{6}k}$, 其中$a_k$定义见上。
    	
    	
    	再具体代入，得到各个模和相位：
    	$a_k\text{的模和相位}$\\
    	$|a_0|=\frac{11-\sqrt{2}}{12},|a_1|=\frac{2+\sqrt{6}}{12}, |a_2|=\frac{2+\sqrt{2}}{12},|a_3|=\frac{1}{12},|a_4|=\frac{2-\sqrt{2}}{12},|a_5|=\frac{\sqrt{6}-2}{12},\\
    	|a_6|=\frac{\sqrt{2}-1}{12},|a_7|=\frac{\sqrt{6}-2}{12},|a_8|=\frac{2-\sqrt{2}}{12},|a_9|=\frac{1}{12},|a_{10}|=\frac{2+\sqrt{2}}{12},|a_{11}|=\frac{2+\sqrt{6}}{12}$\\
    	$a_k\text{的相位均为}0$.
    \end{enumerate}
\end{framed}




\newpage
\section{[40pts]}
求下列信号的离散时间傅里叶变换
\begin{enumerate}[(1)]
    \item $2^{n}\cdot u[-n]$
    \item $(a^{n}\cos{\omega_{0}n})u[n],\left | a \right | < 1$
\end{enumerate}

\begin{framed}
	解：
     \begin{enumerate}[(1)]
    	\item 
    	利用定义即可计算出原信号离散时间傅里叶变换为：
    	\begin{align*}
    		&\quad X(e^{jw}) \\
    		&= \sum\limits_{n = -\infty}^{\infty}x[n]e^{-jwn}\\
    		&= \sum\limits_{n = -\infty}^{\infty}2^n\cdot u[-n] e^{-jwn}\\
    		&= \sum\limits_{n = -\infty}^{0}2^n\cdot e^{-jwn}\\
    		&= \sum\limits_{m = 0}^{\infty}2^{-m}\cdot e^{jwm}\\
    		&= \sum\limits_{m = 0}^{\infty}{(\frac{e^{jw}}{2})}^m\\
    		&= \frac{2}{2-e^{jw}}.
    	\end{align*}
    	
    	\item
    	\begin{align*}
    		&\quad (a^ncos\omega_0n)u[n] \\
    		&= a^n \frac{e^{j\omega_0n} + e^{-j\omega_0n}}{2}u[n]
    	\end{align*}
    	
    	已知单边指数信号的离散傅里叶变换对如下：
    	
    	$a^nu[n]  \leftrightarrow \frac{1}{1-ae^{-j\omega}} \quad(|a| < 1)$.
    	
    	所以利用频移特性可得原信号的离散时间傅里叶变换为：
    	
    	\[\frac{1}{2}(\frac{1}{1-ae^{-j(\omega-\omega_0)}} + \frac{1}{1-ae^{-j(\omega+\omega_0)}}).\]
    	
    \end{enumerate}
\end{framed}





\newpage
\section{[20pts]}
若已知有限长序列$x(n)$如下式
\begin{equation}
x(n) = \begin{cases}
1 & $(n=0)$ \\
2 & $(n=1)$ \\
-1 & $(n=2)$ \\
3 & $(n=3)$
\end{cases}
\end{equation}
求DFT$[x(n)]=X(k)$,再由结果求IDFT$[X(k)]=x(n)$,验证你的计算是正确的（写作矩阵）.

\begin{framed}
   	解：\\
   	由
   	\[
   	X(k) = \sum_{n=0}^{N-1} x(n)W^{kn},
   	\]
   	得矩阵形式：
   	\[
   	\begin{bmatrix}
   		X(0) \\ X(1) \\ X(2) \\ X(3)
   	\end{bmatrix}
   	=
   	\begin{bmatrix}
   		W^0 & W^0 & W^0 & W^0 \\
   		W^0 & W^1 & W^2 & W^3 \\
   		W^0 & W^2 & W^4 & W^6 \\
   		W^0 & W^3 & W^6 & W^9
   	\end{bmatrix}
   	\begin{bmatrix}
   		x(0) \\ x(1) \\ x(2) \\ x(3)
   	\end{bmatrix}.
   	\]
   	
   	由 $W = e^{-j\frac{2\pi}{N}}, N = 4$，得：
   	\[
   	\begin{bmatrix}
   		X(0) \\ X(1) \\ X(2) \\ X(3)
   	\end{bmatrix}
   	=
   	\begin{bmatrix}
   		1 & 1 & 1 & 1 \\
   		1 & -j & -1 & j \\
   		1 & -1 & 1 & -1 \\
   		1 & j & -1 & -j
   	\end{bmatrix}
   	\begin{bmatrix}
   		x(0) \\ x(1) \\ x(2) \\ x(3)
   	\end{bmatrix}
   	=
   	\begin{bmatrix}
   		5 \\ 2 + j \\ -5 \\ 2 - j
   	\end{bmatrix}.
   	\]
   	
   	求 IDFT 类似：
   	\[
   	x(n) =
   	\frac{1}{N}
   	\begin{bmatrix}
   		W^0 & W^0 & W^0 & W^0 \\
   		W^0 & W^{-1} & W^{-2} & W^{-3} \\
   		W^0 & W^{-2} & W^{-4} & W^{-6} \\
   		W^0 & W^{-3} & W^{-6} & W^{-9}
   	\end{bmatrix}
   	\begin{bmatrix}
   		X(0) \\ X(1) \\ X(2) \\ X(3)
   	\end{bmatrix}.
   	\]
   	
   	代入可得：
   	\[
   	x(n) =
   	\frac{1}{4}
   	\begin{bmatrix}
   		1 & 1 & 1 & 1 \\
   		1 & j & -1 & -j \\
   		1 & -1 & 1 & -1 \\
   		1 & -j & -1 & j
   	\end{bmatrix}
   	\begin{bmatrix}
   		5 \\ 2 + j \\ -5 \\ 2 - j
   	\end{bmatrix}
   	=
   	\begin{bmatrix}
   		1 \\ 2 \\ -1 \\ 3
   	\end{bmatrix}.
   	\]
   	
   	与原 $x(n)$ 一致。说明计算正确。
\end{framed}


\end{document} 